% Generated draft: requires mathematical and source review.
\documentclass[11pt]{article}
\usepackage[margin=25mm]{geometry}
\usepackage{fontspec}
\setmainfont{Latin Modern Roman}
\usepackage{amsmath,amssymb,mathtools,unicode-math}
\setmathfont{Latin Modern Math}
\usepackage{enumitem,needspace,xurl,hyperref}
\setlength{\parindent}{0pt}
\setlength{\parskip}{5pt}
\setlength{\emergencystretch}{3em}
\newcommand{\N}{\mathbb{N}}
\newcommand{\Z}{\mathbb{Z}}
\newcommand{\Q}{\mathbb{Q}}
\newcommand{\R}{\mathbb{R}}
\newcommand{\C}{\mathbb{C}}
\newcommand{\F}{\mathbb{F}}
\newcommand{\A}{\mathbb{A}}
\newcommand{\PP}{\mathbb{P}}
\newcommand{\E}{\mathbb{E}}
\newcommand{\eps}{\varepsilon}
\newcommand{\dd}{\,\mathrm{d}}
\begin{document}
\section*{UnsolvedMath ID: 30001138}
\textbf{Source problem code:} catalogue entry
\subsection{Definitions and mathematical statement}
Let a convex four-polytope be a set \(P=\operatorname{conv}(V)\subset\R^4\), where \(V\) is finite and \(\dim P=4\). Its boundary \(\partial P\) is homeomorphic to \(S^3\). Let \(G(P)\) be its one-skeleton, regarded as an abstract graph, and let \(\iota_P:|G(P)|\hookrightarrow\partial P\) be the embedding given by its vertices and edges. Here \(|G(P)|\) denotes the geometric realization of the graph, and a cycle means a simple cycle. Two disjoint cycles have no common vertices.

For a precise conventional interpretation, take embeddings to be piecewise linear and call two disjoint embedded cycles in \(S^3\) \emph{split} if there is a piecewise-linear embedded sphere \(\Sigma\cong S^2\), disjoint from both cycles, such that the cycles lie in different components of \(S^3\setminus\Sigma\). Call the pair \emph{linked} if it is not split. An embedding of a graph is linked if some two disjoint cycles have linked images, and is linkless if every pair of disjoint cycles has split images. These definitions apply in \(\partial P\) through a piecewise-linear homeomorphism with \(S^3\).

Under these conventions, the question is whether there exist a convex four-polytope \(P\) and a piecewise-linear embedding \(e:|G(P)|\hookrightarrow S^3\) such that
\[
\begin{aligned}
&\exists\, C_1,C_2\subset G(P),\quad C_1\cap C_2=\varnothing,
&&\iota_P(C_1)\cup\iota_P(C_2)\text{ is a nonsplit link in }\partial P,\\
&\forall\, D_1,D_2\subset G(P)\text{ that are disjoint cycles},
&&e(D_1)\cup e(D_2)\text{ is a split link in }S^3,
\end{aligned}
\]
where \(C_1,C_2\) in the first line are also cycles.

The supplied source defines a linked graph embedding only by the presence of two disjoint linked cycles. It does not specify the category of embeddings or whether ``linked'' means nonsplit, nonzero linking number, or another condition. Thus the precise formulation above is conditional on the stated piecewise-linear, nonsplit-link interpretation. The ambient space for the alternative embedding is read as \(S^3\) from the context.
\subsection{Short English statement}
Does a convex four-polytope have a linked natural embedding of its edge graph in its boundary, while the same abstract graph admits a linkless embedding in \(S^3\)? The displayed formulation uses piecewise-linear embeddings and interprets linked pairs as nonsplit links; these conventions require source confirmation.
\subsection{Sources}
\begin{itemize}
\item Supplied local registry record: \texttt{data/unsolvedmath/problems.json}, permanent UnsolvedMath ID \(30001138\), title \emph{Linked Skeletons of Convex Four-Polytopes}, exact source problem code \texttt{catalogue entry}; collection: \emph{Oberwolfach Reports Open Problems}. The supplied record gives the year as 2009 and names no individual proposer.
\item Source report: \emph{Discrete Differential Geometry} (2009), \emph{Oberwolfach Reports}, pp.\ 75--144, DOI \texttt{10.4171/owr/2009/02}. \url{https://doi.org/10.4171/owr/2009/02}. The supplied background describes the report as documenting workshop discussions and open problems.
\item Supplied literature reference: Izmestiev, \emph{Colin de Verdiere number and graphs of polytopes}. \url{https://arxiv.org/abs/0704.0349}.
\item Supplied literature reference: Rote, \emph{Open Problems in Discrete Differential Geometry}. \url{https://page.mi.fu-berlin.de/rote/Kram/OWR-DDG09-problems.pdf}.
\item Supplied literature assessment, checked on \texttt{2026-08-21}: the record labels the problem open and states that the specific existence question remains posed and that the checked Colin de Verdiere results do not resolve it. This is supplied dated curation derived from the OWR material, not an independently verified assessment. Time-sensitive status and cited claims require independent verification.
\end{itemize}
\end{document}
